\begin{definition}
Identify $\mathbb{C}$ with $\mathbb{R}^2$ by $z=x+iy$, $(x,y)\in\mathbb{R}^2$.
For a map $f\colon\mathbb{C}\to\mathbb{C}$ let $Df(z)$ denote the real
Fr\'echet derivative of $f$ at $z$, a real-linear map
$\mathbb{C}\to\mathbb{C}$, so that $Df(z)[1]=f_x(z)$ and $Df(z)[i]=f_y(z)$.
The operator of complex differentiation
\[
  \frac{\partial}{\partial z}=\frac12\Bigl(\frac{\partial}{\partial x}-i\frac{\partial}{\partial y}\Bigr)
\]
of Section~2 of the paper is then given by
\[
  \frac{\partial f}{\partial z}(z)=\frac12\bigl(Df(z)[1]-i\,Df(z)[i]\bigr).
\]
We write $f_z$ for this function of $z$, and $H_\xi$ when the variable is
called $\xi$.

At a point where $f$ is not real-differentiable, $Df(z)$ is taken to be the
zero map, so the expression above is defined at every point of $\mathbb{C}$.
\end{definition}
